\section{问题建模}

\subsection{巡测场景与卫星遥感数据集}

东海赤潮应急巡测场景建立在八年真实卫星遥感叶绿素与海表温度大数据之上，航行器遇到的真实海洋由任务当天的实测卫星图构建，而航行器出发时仅携带3天前的旧卫星图。本文使用的数据集源自美国国家海洋和大气管理局海岸观测网的逐日多星融合叶绿素 a 产品，该产品融合了多颗极轨卫星水色传感器的观测，并通过经验正交函数插补算法填补云层遮挡区\cite{beckers2003dineof}，空间分辨率为$2\text{ km}$，时间跨度为2018年1月1日至2026年10月23日共3197天，覆盖东海北纬$25^\circ\text{--}33^\circ$、东经$119.5^\circ\text{--}127^\circ$海域（每日$385\times 361$个网格）。辅助时空预测实验还使用了喷气推进实验室制作的多尺度融合海表温度逐日产品（原始分辨率$0.01^\circ$，重采样至$0.04^\circ$）\cite{chin2017mur}。全文除特别说明外，所有叶绿素 a 浓度、误差标准差与遗憾指标均以对数浓度$\log_{10}(\text{叶绿素 a 浓度}/\text{mg m}^{-3})$为单位。

为构建贴近实战的离岸赤潮应急任务集，本文从每年4月至6月藻华高发期的卫星序列中按统一判据切出正方形巡测窗口。具体而言，在浙江外海陆架区（北纬$26.5^\circ\text{--}31.5^\circ$，东经$121.5^\circ\text{--}125^\circ$）检索离岸至少$30\text{ km}$、经空间平滑后局部峰值超过$5\text{ mg m}^{-3}$的藻华中心，将该中心随机平移至多$16\text{ km}$后截取$80\text{ km}\times 80\text{ km}$（$40\times 40$个$2\text{ km}$网格）的任务窗口，将距陆地或岛屿$15\text{ km}$以内的网格设为禁航区，并筛选可航行水域占比不低于80\%且窗内对数浓度极差不低于0.3的窗口。按此标准共获得422个赤潮任务窗口，其中2018至2022年的274个窗口作为开发集，2023至2026年的148个窗口作为独立的测试集。每次任务开始时，航行器从窗口边缘随机选取的一个可航行网格入水出发。鉴于静止轨道水色卫星逐时观测表明东海表层赤潮斑块主要随表层水体与潮汐整体往复漂移\cite{lou2014diurnal}，本文将航行器运动、叶绿素分布与数据同化建立在随表层水体漂移的参考系中，此时斑块相对于水体保持静止，而海流平流与生消过程的作用则完整体现在3天滞后卫星图与任务当天卫星图之间的实测演变统计中。

为同时检验方法对卫星未分辨小尺度结构的适应性以及在纯实测遥感场上的表现，本文设置两种真实海洋模型。主设置称为“当天卫星图 + 小尺度变化”，在任务当天的真实卫星图上叠加一个标准差为$s_{\text{sat}}=0.15$、空间相关长度为3个网格（$6\text{ km}$）的零均值高斯随机场，用以表征2 km插补卫星图平滑掉的亚像素斑块与反演残差，其幅值依据海洋一号丙星水色传感器在东海与南海现场验证所得约39\%的无偏相对百分比差异（折合对数空间约0.15）标定\cite{ye2021cocts}。第二种设置称为“当天卫星图”，直接以任务当天的真实卫星叶绿素图本身作为真实海洋（即$s_{\text{sat}}=0$）。在两种设置下，航行器出发时携带的先验底图均为任务前3天的真实卫星图$c_{\text{map}}(\mathbf{x})$。

\subsection{航行器动作空间与能耗模型}

机载计算与机动观测的资源冲突由百米级轻型自主水下航行器的机电功耗特性与20分钟离散控制周期共同刻画。参照典型轻型航行器参数\cite{remus100}，航行器对水巡航速度为$1.5\text{ m/s}$（约3节），机载锂电池总能量为$1\text{ kWh}$，在3节巡航状态下续航约22小时，对应巡航总功率$45\text{ W}$，其中基础平台功率为$15\text{ W}$（含导航、荧光计传感器及常驻后台运行的低功耗处理器），推进电机功率为$30\text{ W}$。机载高功耗芯片参照嵌入式边缘图形处理器模块设定（典型功耗范围$10\text{--}25\text{ W}$）\cite{jetsonorinnx}，开启深想时额外消耗$+20\text{ W}$功率（在“省电加速器”硬件配置下额外消耗$+5\text{ W}$），其航线推演速度为低功耗处理器的50倍。由于本文在二维表层场上采用的双航路点折线规划器是对真实水下三维路径评估（需同时积分海流、深度剖面与地形避障约束）的简化替代，决定规划质量的关键因素是机载处理器相对于规划问题的推演速度，因此实验以低功耗处理器每分钟完成的推演次数表征基础算力（取5、20或100次/分，并在5至200次/分范围内开展参数扫描）。

在时长$T=24\text{ h}$的任务窗口内，时间按每段$\Delta t = 20\text{ min}$划分为72个连续决策周期，每个周期内航行器的机体动作与大脑规划并行执行\cite{xiao2020thinking, elboher2023formal}。在机体动作方面，若航行器选择“走”，则开启推进器（$30\text{ W}$）沿当前航线以$1.5\text{ m/s}$航行$1.8\text{ km}$，机载荧光计每航行$1\text{ km}$独立采样一次（每周期获得约2个样本），在电量充足时24小时累计航程约$120\text{--}130\text{ km}$，约为$80\text{ km}$窗口宽度的1.5倍。若航行器选择“停下”，则关停推进器（推进功率为$0\text{ W}$，仅消耗$15\text{ W}$平台功率）随水体漂流，荧光计每4分钟对同一水体采样一次，每周期共获得5个独立样本。荧光计单次测量值$y$与所在网格真实对数浓度$c(\mathbf{x})$的关系满足高斯观测模型：
\begin{equation}
  y = c(\mathbf{x}) + \varepsilon, \quad \varepsilon \sim \mathcal{N}(0, 0.1^2)
  \label{eq:obs}
\end{equation}
在大脑规划方面，每个周期航行器可选择仅维持低功耗处理器后台规划（额外功耗$0\text{ W}$），或开启高功耗芯片“深想”（额外功耗$+20\text{ W}$或$+5\text{ W}$，本周期推演次数提升至后台的50倍）。为消除规划计算耗时带来的位置滞后，新航线始终以航行器在当前20分钟周期结束时刻的预计位置为起点进行搜索，并在周期结束瞬间无缝替换旧航线。当任务满24小时或电池电量耗尽时，航行器停止机动与采样。

\subsection{实时报告机制与时间平均遗憾目标}

应急巡测的目标是尽早找到并持续锁定窗口内叶绿素浓度最高的水体，因此评价指标必须同时奖励最终定位精度与到达高浓水体的时间效率。在$T=24\text{ h}$作业期间的任意时刻$t \in [0, T]$，应急指挥中心均可查询当前已知的最浓水体位置。为防止航行器凭借未经证实的旧图先验妄下结论，规定航行器只能从截至时刻$t$已实际造访过的网格集合$\mathcal{V}(t)$中，报告当前内部集合估计均值$\bar{c}_t(\mathbf{x})$最高的位置$\hat{\mathbf{x}}(t)$：
\begin{equation}
  \hat{\mathbf{x}}(t) = \arg\max_{\mathbf{x}\in\mathcal{V}(t)} \bar{c}_t(\mathbf{x})
  \label{eq:answer}
\end{equation}
若航行器在$24\text{ h}$前因电池耗尽而停机，则停机后的报告位置$\hat{\mathbf{x}}(t)$保持冻结直至$T=24\text{ h}$结束。定义单次任务的时间平均遗憾$R$为全窗口真实最高浓度$\max_{\mathbf{x}} c(\mathbf{x})$与航行器实时报告位置真实浓度$c\big(\hat{\mathbf{x}}(t)\big)$之差在24小时内的时间均值：
\begin{equation}
  R = \frac{1}{T}\int_0^T \left[\max_{\mathbf{x}} c(\mathbf{x}) - c\big(\hat{\mathbf{x}}(t)\big)\right] dt
  \label{eq:regret}
\end{equation}
时间平均遗憾$R$越低表示巡测表现越好，越早抵达真正的高浓度核心区并维持足够的续航时间，积分遗憾就越小。在对数叶绿素尺度下，$R=0.1$对应航行器报告位置的叶绿素浓度比窗口真实最高值低约20\%，而$R=0.3$则意味着报告浓度仅约为真实峰值的二分之一。

\subsection{问题剖析：计算与观测共享有限预算的三重矛盾}

上述巡测模型的本质困难在于内部航线计算与外部机动采样共享同一块$1\text{ kWh}$电池和同一个$24\text{ h}$窗口，而以下三重结构性矛盾直接决定了第4节方法的设计选择。第一，增加计算的真实收益只能在航行器内部的想象海洋中估算，但规划器在搜索样本上的自评得分存在系统性正向偏差。当高功耗芯片将单周期推演次数扩大50倍时，优化器能够精细调整折线航路点以穿过搜索样本集中的偶然高值网格，使内部得分持续攀升但实际航线早已不再改善\cite{smith2006curse, vanhasselt2010double}，这就要求必须将用于搜索航线的想象海洋与用于评判价值的想象海洋彻底分离。

第二，现场观测的信息收益取决于3天前旧卫星图究竟错在何处、错到何种程度，而海洋遥感大数据虽能精确标定这一误差的统计分布，却无法靠数值外推直接消除小尺度误差。由于云区插补平滑与水体亚中尺度演变，确定性或深度学习时空预测均无法从日均旧图中恢复数日后的小尺度峰值位置\cite{assran2025vjepa2}，航行器必须亲赴现场采样，并利用历史卫星大数据标定的误差方差与空间相关尺度构建先验想象海洋，再随每个现场样本通过集合数据同化动态收缩未访水域的不确定性。

第三，最优的时间分配策略随底图不确定性、机载处理器速度及高功耗芯片功耗的组合而发生相变，但单次赤潮任务的天然环境噪声过大，无法靠少量历史任务在线或离线拟合出普适规则。在算力捉襟见肘的慢计算机上，仅靠后台规划会导致航行器在窗口中部盲目绕圈，而在算力充裕的快计算机上，全程开启高功耗芯片只会白白透支推进电池并缩短总航程。由于任务间遗憾标准差远大于策略间平均差异，唯有在下水前针对当前机载硬件参数与数据标定的留出海洋开展全任务周期的完整预演，才能在不依赖任何历史巡测任务的条件下稳健剔除劣质策略。
